% Local dyadic lemma: full abelian scheme, no abelian-quotient cotangent sequence.
\begin{lemma}\label{lem:dyadic-formal-Hecke-kernel}
Let $A/\mathbb Z_2$ be an abelian scheme, and let $T,\delta$ be
commuting endomorphisms over $\mathbb Z_2$. Suppose that their
reductions satisfy
\[
 \overline T=F+\overline\delta V,
\]
where $F,V$ are Frobenius and Verschiebung on $A_{\mathbb F_2}$.
Then $M=T-\delta$ kills the torsion kernel of
$A(\mathbb Q_2)\to A(\mathbb F_2)$.
\end{lemma}
\begin{proof}
Choose integral parameters $t=(t_1,\ldots,t_g)$ at the identity,
and let $A_n$ be the subgroup with $t\in2^n\mathbb Z_2^g$.
The group $A_2$ is torsion free.
Indeed, if a nonzero point has minimum coordinate valuation
$n\geq2$, then $[2](t)=2t+O(t^2)$ has minimum valuation $n+1$.
Multiplication by an odd integer preserves the minimum valuation.
Thus no nonzero point of $A_2$ is torsion.

Let $P\in A_1$ be torsion. Since $[2]A_1\subseteq A_2$, we have
$2P=0$. Write $t(P)=2u$, with $u\in\mathbb Z_2^g$.
Put $H=dV$ and $D=d\overline\delta$, as matrices over $\mathbb F_2$.
The identity $[2]=V\circ F$ on the special fibre shows that the
quadratic part of the multiplication-by-$2$ series, reduced
modulo $2$, is $H t^{[2]}$, where $t^{[2]}=(t_1^2,\ldots,t_g^2)$.
Its linear part over $\mathbb Z_2$ is $2t$.
Consequently, with any integral lift of $H$,
\[
 0=[2](2u)\equiv4u+4H u^{[2]}\pmod8.
\]
Since the residue field is $\mathbb F_2$, this gives
$(H-I)\overline u=0$.

Frobenius has zero differential, so
\[
 d\overline M=D(H-I).
\]
Using the integral series for $M$ now gives
\[
 t(MP)\equiv2\,dM\,u\equiv0\pmod4.
\]
Thus $MP\in A_2$. It is torsion, and hence $MP=0$.
\end{proof}

\begin{corollary}\label{cor:dyadic-Cartier-formal-immersion}
Let $N\geq5$ be odd, $X=X_1(N)$ and $J=J_1(N)$.
On $H^0(X_{\mathbb F_2},\Omega^1)$ put
\[
 U=(\mathcal C-I)H^0(X_{\mathbb F_2},\Omega^1),
\]
where $\mathcal C$ is Cartier.
Let $E\in X^{(d)}(\mathbb F_2)$.
If restriction from $U$ to $H^0(E,\Omega^1_X|_E)$ is surjective,
then $(T_2-\langle2\rangle)\circ\iota$ is a formal immersion at $E$,
where $\iota:X^{(d)}\to J$ is an integral Abel--Jacobi map.
In particular, if two points $D,C\in X^{(d)}(\mathbb Q_2)$ reduce
to $E$ and $[D-C]\in J(\mathbb Q_2)$ is torsion, then $D=C$.
\end{corollary}
\begin{proof}
The full Jacobian has good reduction at $2$.
Eichler--Shimura gives $\overline T_2=F+\langle2\rangle V$.
Under the Abel--Jacobi identification of invariant differentials
with $H^0(X_{\mathbb F_2},\Omega^1)$, pullback by $V$ is Cartier.
Therefore, for $M=T_2-\langle2\rangle$ and $\delta=\langle2\rangle$,
\[
 M^*=(\mathcal C-I)\circ\delta^*.
\]
Since $\delta^*$ is invertible, $\operatorname{im}(M^*)=U$.

The cotangent map of the Abel--Jacobi morphism at $E$ is the
restriction map to $H^0(E,\Omega_X^1|_E)$.
This description holds for nonreduced divisors and nonrational
support points: it follows from the perfect residue pairing with
$H^0(E,\mathcal O_E(E))$, the tangent space of $X^{(d)}$ at $E$.
The assumed surjectivity therefore gives a formal immersion.
This uses the endomorphism of the full smooth Jacobian directly;
no cotangent sequence for an abelian quotient is required.

For the final assertion, properness extends $D,C$ over
$\mathbb Z_2$. Their torsion difference reduces to zero, so
Lemma~\ref{lem:dyadic-formal-Hecke-kernel} gives $M[D-C]=0$.
Their images under $M\circ\iota$ are equal.
Formal immersion is injective on the corresponding residue class,
since its map of completed local rings is surjective.
Hence $D=C$.
\end{proof}

\noindent\emph{Local cuspidal traces.}
For an application with $D$ defined over $\mathbb Q$ and $J(\mathbb Q)$
of rank zero, the cuspidal lift $C$ need only be defined over
$\mathbb Q_2$.
The cuspidal scheme is finite etale at $2$, so a closed residue cusp
of degree $f$ lifts to an unramified degree-$f$ cusp over
$\mathbb Q_2$; its trace is a $\mathbb Q_2$-rational divisor.
Manin--Drinfeld makes $[C-dc_0]$ torsion, while rank zero makes
$[D-dc_0]$ torsion, for a rational cusp $c_0$.
Thus $[D-C]\in J(\mathbb Q_2)$ is torsion, even when individual
cusps in $C$ have degree $3$.
The residue field used in the kernel lemma remains $\mathbb F_2$.
